\documentclass[10pt,twocolumn]{article}

\usepackage[top=1in,bottom=1in,left=0.85in,right=0.85in,columnsep=0.35in]{geometry}
\usepackage{amsmath}
\usepackage{amssymb}
\usepackage{stmaryrd}
\usepackage{booktabs}
\usepackage{microtype}
\usepackage{titlesec}
\usepackage{fancyhdr}
\usepackage{hyperref}
\usepackage{listings}
\usepackage{xcolor}
\usepackage{times}
\usepackage{setspace}
\usepackage{flushend}

\setlength{\parskip}{3pt plus 1pt minus 1pt}
\setlength{\parindent}{1em}
\titleformat{\section}
  {\normalsize\bfseries\scshape}{}{0em}{}
\titlespacing*{\section}{0pt}{2.2ex plus 0.6ex minus 0.3ex}{0.5ex plus 0.1ex}
\titleformat{\subsection}
  {\normalsize\itshape}{}{0em}{}
\titlespacing*{\subsection}{0pt}{1.4ex plus 0.4ex minus 0.2ex}{0.3ex}
\pagestyle{fancy}
\fancyhf{}
\renewcommand{\headrulewidth}{0pt}
\fancyhead[C]{\small\textit{sh2elf: Compiling POSIX Shell to Native x86-64}}
\fancyfoot[C]{\small\thepage}

\fancypagestyle{firstpage}{%
  \fancyhf{}
  \renewcommand{\headrulewidth}{0pt}
  \fancyfoot[C]{\small\thepage}
}

\lstset{
  basicstyle=\ttfamily\footnotesize,
  breaklines=true,
  columns=fullflexible,
  keepspaces=true,
  frame=none,
  aboveskip=3pt,
  belowskip=3pt,
}

\hypersetup{
  colorlinks=true,
  linkcolor=black,
  urlcolor=black,
  citecolor=black,
  pdftitle={sh2elf: Compiling POSIX Shell Scripts to Native x86-64 ELF Binaries},
  pdfauthor={Ivan Gaydardzhiev},
}

\begin{document}

\thispagestyle{firstpage}

\twocolumn[{%
\vspace*{10pt}
\begin{center}
  {\LARGE\bfseries sh2elf: Compiling POSIX Shell Scripts\\[5pt]
  to Native x86-64 ELF Binaries}\\[12pt]
  {\large Ivan Gaydardzhiev}\\[3pt]
  {\normalsize\itshape Independent Research}\\[10pt]
  {\small This work is licensed under the Creative Commons
  Attribution 4.0 International License (CC~BY~4.0).}\\[18pt]
\end{center}

\begin{center}
\begin{minipage}{0.82\textwidth}
\small\noindent\textbf{Abstract.}\enspace
sh2elf is a single-pass compiler that translates POSIX shell scripts
into statically linked ELF64 executables for Linux on x86-64. It
handles the full shell language subset including pipelines, control
flow, parameter expansion, arithmetic, signal handling, and a large
set of built-in commands, all without invoking an external shell at
runtime. The compiler is written in C, totals roughly 3100 lines, and
produces binaries that execute 11.28~times faster than the equivalent
bash script on a comprehensive benchmark exercising the entire
supported language. The output binary carries no shared library
dependencies and communicates with the kernel exclusively through raw
Linux system calls.
\end{minipage}
\end{center}
\vspace{18pt}
}]


\section{Introduction}

The POSIX shell is not a fast language. Every invocation parses text,
forks subprocesses to evaluate command substitutions, and interprets
variable expansions at runtime. For scripts that run on hot paths or
that are shipped as part of a distribution's init or package
infrastructure, this overhead is constant and unavoidable. Attempts to
speed up shell execution have historically fallen into two categories:
faster interpreters such as dash, and transpilation to another
interpreted language such as Python or Perl. Neither approach
eliminates the interpreter from the execution path.

sh2elf takes a different position. Given a shell script, it compiles
it once and produces a native ELF64 binary that the Linux kernel can
execute directly. The resulting program contains no interpreter, no
shared libraries, and no calls back into the shell. All control flow
is native branches. All I/O goes through inline system call sequences.
Variable state lives in a BSS hashtable. Pipelines are implemented
with \texttt{fork}, \texttt{pipe}, and \texttt{dup2} at the call site,
emitted as machine code.

The practical consequence is a binary that runs an order of magnitude
faster than bash on identical workloads and that can be copied to any
x86-64 Linux system and executed without any runtime dependencies.


\section{Compiler Architecture}

The compiler is structured as a front end, an intermediate
representation, a code generator, and an ELF writer. All four stages
are contained in \texttt{sh2elf.c}, a single C source file. There are
no external parser generators, no LLVM passes, and no linking to a C
runtime.

\subsection{Front end and IR}

The front end reads a shell script and produces a flat array of IR
statements. Each IR statement holds a condition field indicating
whether it participates in a \texttt{\&\&} or \texttt{||} chain, and a
pipeline field that lists the stages of the pipeline as argument
vectors. This is a three-address representation in spirit: the front
end resolves variable references and quoting rules into slot indices
rather than keeping raw text around, though for dynamic values the
full resolution is deferred to runtime code sequences.

The IR layer decouples parsing from code generation cleanly enough
that the compiler can also emit a textual dump of the IR for
inspection, and can produce a source-annotated disassembly listing
that maps each shell line to the x86-64 instructions it generated.

\subsection{Code generator}

The code generator walks the IR and emits x86-64 machine code
directly into a growable byte buffer. There is no assembler text
involved at any stage. Each instruction is written byte by byte using
small helper functions that encode register-immediate and
register-register forms.

For built-in commands, the generator emits inline system call
sequences. A \texttt{pwd} statement becomes a
\texttt{mov~rdi} / \texttt{mov~rsi} / \texttt{mov~rax,~79} /
\texttt{syscall} sequence. A \texttt{printf} with a known string
argument becomes a \texttt{write} syscall with the string address
patched in from the string pool. For external commands, the generator
emits a \texttt{fork}/\texttt{execve} pair at the call site. The
compiler selects the full path from a fixed list of prefixes at
compile time, so the runtime does not search \texttt{PATH}.

Arithmetic expansion is handled by a 64-bit recursive descent integer
evaluator that emits x86-64 arithmetic instructions directly. The
expression \texttt{\$(( A * B + C ))} compiles to a small sequence of
\texttt{mov}, \texttt{imul}, and \texttt{add} instructions operating
on \texttt{rax}.

\subsection{String pool}

All string literals are collected into a \texttt{StrPool} at compile
time. The pool uses linear search for deduplication: before adding a
string it scans for an existing match, and duplicate strings are not
emitted. The pool is placed in the binary in a read-only segment, and
the code generator emits relocations that the ELF writer resolves
before writing the file.

\subsection{ELF writer}

The ELF writer constructs the output binary entirely in memory and
writes it in a single call. It builds an ELF64 header with two program
headers: one \texttt{PT\_LOAD} segment covering the text and
read-only data at virtual address \texttt{0x400000} with read and
execute permissions, and one \texttt{PT\_LOAD} segment covering the
BSS at \texttt{0x600000} with read and write permissions. The entry
point is set to the first generated instruction, which immediately
follows the ELF and program headers in the file.

Before writing, the ELF writer runs the peephole optimizer on the
code buffer. The optimizer replaces \texttt{mov~rax,~0} with
\texttt{xor~eax,~eax} and removes unconditional jumps that target the
following instruction. It then patches all pending relocations,
filling in the absolute virtual addresses for string pool entries that
the code generator left as placeholders.


\section{Memory Model}

The generated binary uses three memory regions. The text segment holds
executable code beginning at \texttt{0x400000}. Immediately following
the code in the same segment is the read-only string pool. The BSS
segment begins at \texttt{0x600000} and holds all mutable state.

The BSS layout is fixed across all generated binaries. Offset~0 holds
the exit status register, an 8-byte slot that every command updates
with its return code. Offset~8 holds the positional parameter frame:
slots for \texttt{\$0} through \texttt{\$9}, plus \texttt{\$\#},
\texttt{\$@}, and \texttt{\$*}. These are populated by the binary
prologue, which reads \texttt{argc} and \texttt{argv} from the Linux
initial stack before executing any shell statements.

Offset~256 begins the variable hashtable, which uses open addressing
with linear probing. Each slot holds a key pointer and a value
pointer. Variable lookup and assignment are emitted as inline code
sequences that hash the key string, probe the table, and read or write
the value pointer without any function call overhead, since the table
address is a compile-time constant.

Offset~4352 is the I/O scratch buffer, a 64~KB region used by
streaming built-in commands for their \texttt{read} and \texttt{write}
syscall loops.


\section{Pipeline Implementation}

Pipelines are compiled into a sequence of \texttt{pipe},
\texttt{fork}, \texttt{dup2}, \texttt{close}, and \texttt{wait4}
system calls emitted inline at the point where the pipeline appears in
the source. For a two-stage pipeline \texttt{A~|~B}, the generated
code calls \texttt{sys\_pipe} to create a file descriptor pair, forks
a child for \texttt{A} that redirects its stdout to the write end of
the pipe, and forks a child for \texttt{B} that redirects its stdin to
the read end. The parent closes both ends and waits for both children.
The exit status of the pipeline is taken from the last stage and
stored into the BSS status register.

Pipelines of arbitrary length are supported. Each additional stage
introduces one more \texttt{fork}/\texttt{dup2}/\texttt{close} block,
and the parent accumulates \texttt{wait4} calls for all children,
discarding all statuses except the last.

\texttt{SIGPIPE} handling is set up by a \texttt{sys\_rt\_sigaction}
call in the binary prologue with \texttt{SIG\_DFL}, restoring default
signal disposition for pipeline stages that may receive a broken pipe.


\section{Variable Handling}

Variables are dynamic at runtime. The compiler cannot statically
determine all variable names or their lifetimes in the general case
because of \texttt{eval} and because variables may be set by external
commands through environment inheritance. The BSS hashtable
accommodates this: any variable name encountered at runtime is hashed
and stored as a key-value pair of \texttt{char*} pointers.

For variables whose names and values are known at compile time, the
compiler can sometimes emit the value directly as a string pool
reference and bypass the hashtable entirely. This applies to simple
cases where no variable is involved, and to arithmetic expressions
where both operands are numeric literals.

Parameter expansions are emitted as inline runtime routines. The
expansion \texttt{\$\{VAR\#\#pat\}} compiles to a sequence that scans
the value of \texttt{VAR} from left to right evaluating
\texttt{fnmatch} at each offset, takes the longest prefix that matches
\texttt{pat}, and returns the remaining suffix. The expansion
\texttt{\$\{VAR/pat/rep\}} performs substitution rather than stripping
by the same mechanism. All these operations work directly on string
data in the BSS scratch buffer.


\section{Control Flow}

Control flow constructs compile to native branch sequences with
forward or backward jumps patched after the target address is known.

An \texttt{if} statement compiles to: evaluate the condition pipeline,
load the status from the BSS status register, compare against zero,
and branch forward over the then-body if non-zero. The else branch, if
present, is prefixed by an unconditional jump over it from the end of
the then-body. Patch addresses are recorded on a compile-time stack
and filled in once the target offset is known.

\texttt{while} and \texttt{until} loops record the loop top address
before emitting the condition, emit the condition, branch past the
loop body on failure (or success for \texttt{until}), emit the body,
and emit an unconditional backward jump to the loop top. \texttt{break
N} and \texttt{continue N} are implemented with a backpatching stack
that records patch sites for each nesting level.

\texttt{case} statements compile to a chain of \texttt{fnmatch}
evaluations. Each pattern compiles to a test-and-branch sequence. On a
match the compiler emits the corresponding body followed by an
unconditional jump to the end of the entire \texttt{case} block.


\section{Built-in Commands}

sh2elf implements a large set of shell built-in commands as inline
code rather than external binaries, using raw system calls without any
C standard library involvement. This is the second main source of the
compiler's performance advantage: where bash forks a subprocess to run
\texttt{grep} or \texttt{sort}, sh2elf emits an inline loop of
\texttt{read}/\texttt{write} syscalls that processes the stream
without forking.

The streaming built-ins, including \texttt{grep}, \texttt{sort},
\texttt{uniq}, \texttt{cut}, \texttt{tr}, \texttt{head}, \texttt{tail},
\texttt{sed}, \texttt{awk}, \texttt{find}, \texttt{ls}, \texttt{ps},
\texttt{tar}, \texttt{gzip}, and \texttt{gunzip}, read from stdin
through the BSS I/O buffer and write to stdout through the same
buffer. The 64~KB buffer size matches typical kernel pipe buffer sizes
and keeps memory pressure low.

File-manipulation built-ins such as \texttt{cp}, \texttt{mv},
\texttt{rm}, \texttt{mkdir}, \texttt{rmdir}, \texttt{touch},
\texttt{chmod}, \texttt{chown}, and \texttt{chgrp} are compiled to the
corresponding syscalls directly. Process-related built-ins including
\texttt{kill}, \texttt{killall}, \texttt{pgrep}, \texttt{pkill},
\texttt{ps}, and \texttt{nice} operate by scanning \texttt{/proc}
through \texttt{sys\_getdents64} syscalls, matching process names
against the argument, and issuing \texttt{sys\_kill} on matching PIDs.

The \texttt{eval} built-in is implemented as a runtime compilation
step. The argument string is assembled, written to a temporary file,
and the compiler is reinvoked on that file. This makes \texttt{eval}
correct but not fast; a script that calls it in a hot loop will not
benefit from compilation.


\section{Tooling}

The compiler binary provides several operating modes beyond basic
compilation. The \texttt{{-}{-}list} mode emits a source-annotated
listing that pairs each shell source line with the x86-64 instructions
it generated, useful for verifying code generation correctness. The
\texttt{{-}{-}dump-ir} mode prints the three-address IR before code
generation, showing each statement's condition flag and pipeline
argument vectors. The \texttt{{-}{-}emit-obj} mode produces a
relocatable ELF \texttt{.o} object file that can be linked into a C
program, allowing shell logic compiled by sh2elf to be called from C.

A language server, \texttt{sh2elf-lsp}, implements a subset of the
Language Server Protocol and provides real-time parse error reporting
with line and column positions, and hover information for built-in
commands.

A differential test harness, \texttt{test\_diff}, compiles each test
script with sh2elf and runs both the compiled binary and bash on the
same input, comparing stdout and exit status byte-for-byte. The test
suite covers 89 integration tests across the full language subset, all
of which pass. A grammar-based fuzzer, \texttt{fuzz\_sh}, generates
random but syntactically valid shell scripts and runs them through the
compiler and code generator; 100 iterations completed without
triggering a crash or assertion failure.


\section{Performance}

The benchmark runs a 130-line shell script that exercises every
supported language construct: variables, arithmetic, control flow,
pipelines, parameter expansions, string operations, file manipulation,
and built-in commands. The script is run 50 times for each executor.

\begin{table}[h]
\centering
\small
\setlength{\tabcolsep}{8pt}
\begin{tabular}{lrrr}
\toprule
Executor & Total (50 runs) & Per run & Speedup \\
\midrule
bash     & 3.668\,s & 73.3\,ms & 1.0$\times$ \\
sh2elf   & 0.325\,s &  6.5\,ms & \textbf{11.28$\times$} \\
\bottomrule
\end{tabular}
\caption{Execution time on a comprehensive POSIX benchmark.}
\end{table}

The speedup comes from two sources. First, sh2elf eliminates
interpreter startup and text parsing on every run. Second, built-in
commands that bash implements by forking an external process are
emitted as inline syscall loops in the sh2elf binary, removing the
\texttt{fork}/\texttt{exec}/\texttt{wait} overhead for each
invocation.


\section{Formal Semantics}

The compilation function $\mathcal{C}$ maps a shell script to an ELF64
binary represented as a tuple
$\langle \text{Code},\, \text{StrPool},\, \text{Relocations},\,
\text{BSS} \rangle$.

For simple built-in commands with statically known arguments,
compilation is fully determined at compile time. The command
\texttt{printf~s} compiles to the sequence:
\begin{align*}
\mathcal{C}\llbracket \texttt{printf~}s \rrbracket \;=\;\;
  &\texttt{mov~rdi,\,1} \\
  &\texttt{mov~rsi,\,Addr}(s) \\
  &\texttt{mov~rdx,\,Len}(s) \\
  &\texttt{mov~rax,\,1} \\
  &\texttt{syscall}
\end{align*}

A two-stage pipeline $C_1 \mid C_2$ compiles to a fixed sequence:
\texttt{sys\_pipe} creates the descriptor pair; a first
\texttt{sys\_fork} produces a child that calls
$\texttt{sys\_dup2}(fd_1, 1)$, closes $fd_0$, and executes $C_1$; a
second \texttt{sys\_fork} produces a child that calls
$\texttt{sys\_dup2}(fd_0, 0)$, closes $fd_1$, and executes $C_2$; the
parent closes both descriptors, issues \texttt{sys\_wait4} for both
children, and stores the status of $C_2$ into $\mathrm{BSS}[0]$.

Short-circuit operators compile to conditional branches on the BSS
status register. The \texttt{\&\&} operator compiles as follows:
emit $\mathcal{C}\llbracket P_1 \rrbracket$, then
\texttt{cmp~[BSS\textsubscript{0}],\,0}, then
\texttt{jne~}$L_\text{skip}$, then
$\mathcal{C}\llbracket P_2 \rrbracket$, then place $L_\text{skip}$.
The \texttt{||} operator is symmetric with \texttt{je} in place of
\texttt{jne}.

Parameter stripping expansions of the form \texttt{\$\{VAR\#\#pat\}}
evaluate $\texttt{fnmatch}(pat,\, S[0..k])$ for $k = 1,\ldots,|S|$
and return the suffix after the longest matching prefix. Suffix
stripping \texttt{\$\{VAR\%\%pat\}} mirrors this from the right.
Dynamic substitution \texttt{\$\{VAR//pat/rep\}} scans the buffer and
replaces every non-overlapping match of \texttt{pat} with \texttt{rep}.


\section{Limitations}

sh2elf targets a well-defined subset of POSIX shell. Constructs outside
that subset are rejected with a parse error at compile time rather
than silently miscompiled. The source language deliberately excludes
features that cannot be compiled efficiently or that require a runtime
interpreter: dynamic \texttt{source} of unknown files, co-processes,
job control signals, and interactive features such as readline history
are not supported.

The \texttt{eval} built-in is an exception. Because its argument is an
arbitrary string assembled at runtime, it requires a runtime
compilation step, making it correct but not fast.

The output binary carries no dynamic relocation information. The fixed
virtual addresses are suitable for statically linked executables where
no dynamic linker is involved, which is exactly the use case sh2elf
targets.


\section{Conclusion}

sh2elf demonstrates that the semantic gap between POSIX shell and
native machine code is smaller than it might appear. The core shell
execution model, pipelines, redirections, variable expansion, and
control flow, maps directly to sequences of Linux system calls and
x86-64 branch instructions. A compiler that performs this translation
statically removes the interpreter overhead entirely and produces
binaries that run at speeds appropriate for systems software rather
than scripting.

The compiler is self-contained, requiring no external tools beyond a C
compiler to build. All built-in command implementations use raw system
calls. The resulting binaries have no shared library dependencies and
no runtime requirements beyond the Linux kernel ABI. For shell scripts
that need to be fast, distributable as single files, or embedded in
environments where a shell may not be present, sh2elf provides a
practical compilation path.

\bigskip
\noindent\rule{\columnwidth}{0.4pt}
\smallskip

\noindent{\small
The source code for sh2elf is available under the GNU General Public
License version~3.
Copyright \copyright\ 2025--2026 Ivan Gaydardzhiev.
This paper is licensed under CC~BY~4.0.}

\end{document}
